\section{Pattern extensions: domain ontologies and graph-as-research-infrastructure}
\label{sec:appendix-pattern-extensions}

This appendix expands two forward-looking subsections that the
\paperref{sec:pattern}{pattern section of the main paper} points at: how the pattern composes
with domain-specific ontologies, and how the bioinformatics
precedent for semantic-tagging-as-norm reframes provenance graphs
as research infrastructure under volume pressure. Neither is on
the critical path of the eight integrated practices, but each is
how F(AI)\textsuperscript{2}R generalises beyond the present case.

\subsection{Extensibility through domain ontologies}
\label{sec:domain-ontologies}

PROV-O is the spine; it does not have to be the only ontology in
the graph. The schema we have introduced is compositional: any
\texttt{fair2r:Claim} can carry additional triples that ground its
\emph{values} in a domain ontology, alongside the
\texttt{prov:wasDerivedFrom} edge that grounds its
\emph{evidence}. The two axes are orthogonal --- provenance asks
\emph{where did this claim come from?}, domain ontologies ask
\emph{what does the value of this claim mean?} --- and a manuscript
that answers both is more reproducible than one that answers only
the first.

A neighbouring vocabulary worth naming explicitly is MLCommons
Croissant~1.0~\cite{benjelloun2024croissant}, the schema.org /
JSON-LD vocabulary for ML datasets and their responsible-AI
annotations. Croissant occupies the dataset-metadata slot
(see the \paperref{sec:limits-and-objections}{limits and objections section of the main paper}) and composes cleanly with the
F(AI)\textsuperscript{2}R schema: a \texttt{fair2r:Source} can
reference a Croissant record so that the provenance graph for a
manuscript and the dataset metadata for the data it sits on top of
read as a single linked-data graph rather than two disconnected
ones.

Two examples concrete enough to plug in today:

\begin{description}[leftmargin=1.2em,itemsep=2pt,topsep=4pt]
    \item[\textbf{Sensor data: SOSA / SSN.}] The W3C Semantic Sensor
        Network ontology and its lightweight core
        SOSA~\cite{janowicz2019sosa}\todo[inline]{verify} expose
        \texttt{sosa:Observation}, \texttt{sosa:Sensor},
        \texttt{sosa:FeatureOfInterest}, and the relations between
        them. A \texttt{fair2r:Claim} that asserts a measurement can
        carry a \texttt{prov:wasDerivedFrom} pointing at the
        underlying \texttt{sosa:Observation}, and the audit pass
        gains the ability to ask \emph{is this measurement well-
        formed?} as a SHACL query rather than a paragraph of
        prose.
    \item[\textbf{Units of measure: OM-2 / QUDT.}] The OM-2 ontology
        of units of measure~\cite{rijgersberg2013om}\todo[inline]{verify}
        and the QUDT vocabulary~\cite{qudt}\todo[inline]{verify}
        expose \texttt{om:Quantity}, \texttt{om:Unit},
        \texttt{qudt:hasUnit}, and the dimensional relations between
        units. A claim that asserts a numeric value can carry the
        unit as an explicit triple rather than as a typographic
        decoration. Unit consistency is then a graph query, not a
        manuscript-reading exercise.
\end{description}

\noindent The class of failure modes this addresses is the one
prose alone cannot reach: a value that is internally coherent in
its prose and its citation but \emph{wrong in its units}, or
\emph{wrong in the dimensional sense it implies}. A reviewer
catching such a defect today reads the figure caption and the
methods section and triangulates; an audit pass with a domain
ontology reads the graph and the SHACL shapes and reports.

The recommendation is open. F(AI)\textsuperscript{2}R does not
prescribe specific domain ontologies because the right choice is
discipline-specific. It prescribes that, where a discipline already
has a \emph{konkrete standardisierte Methode} expressed as an
ontology, the F(AI)\textsuperscript{2}R repository should adopt it
and link from \texttt{fair2r:Claim} entities to its terms. NFDI4Ing
and the Helmholtz Metadata Collaboration are the nearest two
communities to this paper that maintain such ontologies for
engineering and metadata respectively, and the integration cost is
low. The ceiling is the discipline's; the floor is what
F(AI)\textsuperscript{2}R asks for.

\subsection{Looking forward: provenance graphs as research infrastructure}
\label{sec:research-infrastructure}

The argument of \S\ref{sec:domain-ontologies} generalises. The
discipline that pioneered semantic-tagging-as-norm under
data-volume pressure is bioinformatics: faced with what was, at the
time, an obscene amount of sequence data, the field standardised
identifier and annotation schemas (the Gene
Ontology~\cite{ashburner2000go}\todo[inline]{verify} being the most
visible example) so that genes, proteins, and their relations could
be referenced unambiguously across labs, papers, and databases. The
ontology was not a research output; it was a precondition for
research that scaled. The community's investment in semantic
tagging was, in retrospect, what made the next decade of
bioinformatics possible.

We claim a closer analogy than is sometimes recognised. LLM-assisted
scholarly writing is producing artefacts at a volume that the
unaided reviewer pipeline cannot absorb (see the \paperref{sec:intro}{introduction of the main paper}). The
response will, on present trajectory, look more like the
bioinformatics response than like the editorial response: not new
detection tools but new shared schemas, ontologies, and
graph-interconnected metadata that let claims, sources, and
activities reference each other across institutional boundaries.
\textbf{Future research success may come to depend less on the
solitary craft of a single manuscript and more on three things the
F(AI)\textsuperscript{2}R schema is designed to support:} (i)
provenance models that turn published claims into machine-readable
hypotheses for the next study; (ii) graph interconnectivity across
institutional knowledge graphs --- the Helmholtz cross-centre
graph~\cite{curdt2025hmc} is the proximate example, and a
fully-instrumented version would let a claim in one centre's
manuscript be linked, queried, and challenged by claims in another's
without leaving the linked-data layer; (iii) clever experimental
validation that uses the graph to identify the next testable
deviation from the published consensus.

That trajectory is a position the paper is willing to take and a
direction the eight integrated practices serve. The verification
ladder feeds (i); the per-claim provenance map and the
domain-ontology side-arms feed (ii); the recursive case study and
the audit pass feed (iii). None of this requires F(AI)\textsuperscript{2}R
to be the standard. It requires only that some standard with these
properties be adopted, and the present paper is one offer among
several the community can refine.
